\section{Level 2: Future Prediction}
\label{sec:level2}

\subsection{Capability Definition and Adjacent Boundary}

Level~2 predicts a future state, observation, event, or trajectory from what has already been observed. It covers longitudinal imaging, disease-progression modeling, EHR event generation, physiological forecasting, and missing-visit synthesis. Its defining question is what is likely to happen next. A model can be continuous-time, probabilistic, autoregressive, or generative and still remain at this level when no settable action enters the transition. Elapsed time, visit index, disease stage, and observed care history are covariates rather than interventions.

The lower boundary is a directed temporal target. A same-time reconstruction remains Level~1 even if its representation is predictive. The upper boundary is a settable action. To reach Level~3, the model must allow an action to be changed for the same pre-action state, use that action inside the transition, and evaluate the resulting post-action future. Predicting the medication a clinician is likely to prescribe is future prediction; simulating a patient's subsequent state after a specified medication can be action-conditioned simulation. This distinction prevents care-process forecasting from being read as intervention response.

Time-conditioned simulation is an instructive boundary. MRI CEKWorld models contrast-enhancement kinetics through $\hat I_p(t)=M_\theta(I_{p,0},t)$ \citep{Kong2026MRIContrastEnhancement}. The model has a clinically meaningful temporal transition and can generate observations at selected times. The selected coordinate is not a treatment or acquisition action, however, so the implemented claim is Level~2. It would reach Level~3 if dose, contrast agent, or acquisition policy were settable inputs whose effects on the future observation were evaluated. The distinction does not diminish the value of continuous-time modeling; it states accurately what question the model answers.

\subsection{Clinical Evidence Landscape}

\subsubsection{Longitudinal Imaging and Disease Progression}

Longitudinal imaging is a central Level~2 application because disease is often observed through repeated scans. Early work learned aging trajectories, disease states, or coupled forecasts of clinical scores and three-dimensional MRI; subsequent generative models introduced explicit temporal embeddings, sequence-aware diffusion, and conditional image synthesis \citep{Xia2021Learningsynthesiseageing,Severson2021DiscoveryParkinsons,Zhao2022Multiviewprediction,Sch_n2023ExplicitTemporalEmbedding,Yoon2023SADMSequenceaware,Dao2024ConditionalDiffusionModel}. This lineage establishes a useful distinction between cross-sectional synthesis and a patient-linked temporal target. It also exposes a recurring limitation: a model trained to match population-level change can produce a plausible future without identifying the individual mechanism that generated it.

Recent systems broaden this line from brain aging and neurodegeneration to tumor growth, pulmonary fibrosis, ophthalmology, and general four-dimensional reconstruction. Latent diffusion, flow matching, continuous-time reconstruction, and multiscale trajectory models represent irregularly sampled evolution at a selected horizon \citep{Puglisi2025BrainLatentProgression,Aghaei2025integratedpredictivemodel,Laslo2025MechanisticLearningGuided,Chen2025LearningPatientSpecific,Kapoor2025MrextrapLongitudinalAging,Zhao20254dVqGan,Disch2025CronosContinuousTime,Disch2025TemporalFlowMatching,Liu2025ImageflownetForecastingMultiscale}. A progression-map-guided model, for example, predicts patient-linked structural MRI at later follow-up in Alzheimer disease, while transition-based digital twins address sparse visits \citep{Huang2026TransitionBasedDigital,Zhao2026ForecastingAD}. Their shared contribution is temporal continuity under incomplete observation. Their shared evidentiary burden is to show that identity, pathology, and uncertainty remain credible as the forecast horizon increases.

Dynamic medical observations extend beyond conventional follow-up scans. Video and motion models synthesize endoscopy, echocardiography, ultrasound, cardiac motion, and general medical video; related systems model brain activity, ocular dynamics, inertial kinematics, and contrast-enhancement kinetics \citep{Cao2024MedicalVideoGeneration,Li2024EndoraVideoGeneration,Li2024EchopulseEcgControlled,Yu2024ExplainableControllableMotion,Caro2024BrainlmFoundationModel,Chen2025UltrasoundImageVideo,You2025TemporalDifferentialFields,Gao2026EyeworldGenerativeWorld,Delaney2026SonataHybridWorld,Kong2026MRIContrastEnhancement}. These studies make the observation process itself dynamic. ECG, motion curves, or time coordinates may control what is generated, but they do not necessarily constitute a clinical intervention. Likewise, pre-to-post-treatment image generation and latent editing can remain Level~2 when treatment is fixed by the dataset rather than exposed as a settable, validated action \citep{Yeganeh2025LatentDriftingDiffusion,Leclercq2025PrePostTreatment}.

\subsubsection{EHR Events and Patient-Trajectory Models}

EHR trajectory modeling developed from patient representations and next-event prediction toward generative foundation models of clinical timelines. Deep Patient and Doctor AI established representation learning and recurrent event forecasting; BEHRT, CEHR-BERT, Hi-BEHRT, and CLMBR then showed how temporal structure and large-scale pretraining could support downstream prediction \citep{Miotto2016DeepPatientUnsupervised,Choi2016DoctorAiPredicting,Li2020BEHRTTransformerElectronic,Pang2021CehrBertIncorporating,Li2021HiBehrtHierarchical,Steinberg2021LanguageModelsAre}. These works are important predecessors even when the learned representation is not itself a world model. They supply the state and temporal abstractions on which later generative trajectories depend.

Encoder--decoder and autoregressive systems subsequently made the future record an explicit modeling target. TransformEHR, MOTOR, Event Stream GPT, Foresight, Foresight~2, zero-shot trajectory prediction, and optimized BEHRT variants cover event sequences, time-to-event structure, and transfer across sites or tasks \citep{Yang2023TransformEHRtransformerbased,Steinberg2023MOTORTimeEvent,McDermott2023EventStreamGpt,Kraljevi_2024Foresightgenerativepretrained,Renc2024ZeroShotHealth,Odgaard2024CoreBehrtCarefully,Kraljevic2024LargeLanguageModels,Guo2024multicenterstudy}. More recent scaling studies extend this family to generative medical events, adaptive risk estimation, multi-task modeling, long-context evaluation, health-system-scale records, multimodal pathways, and recurrence-aware next-visit prediction \citep{Waxler2025GenerativeMedicalEvent,Renc2025Foundationmodelelectronic,Redekop2025Zeroshotmedical,Pang2025CEHRXGPTScalable,Kim2025MultimodalElectronicHealth,Wornow2025ContextCluesEvaluating,Rajamohan2025FoundationModelsClinical,Pellegrini2025Ehr2pathScalableModeling,Rajamohan2026ScalingRecurrenceAware}. Scale improves coverage of clinical sequences, but it does not by itself separate latent patient dynamics from documentation and care delivery.

Synthetic-EHR models form a related branch. Mixed-type longitudinal generators, diffusion models, hierarchical autoregressive systems, and semantics-guided transformers aim to preserve temporal and clinical structure while producing reusable trajectories \citep{Li2023Generatingsyntheticmixed,Kuo2023SyntheticHealthrelated,Theodorou2023SynthesizeHighDimensional,Zhou2025GeneratingClinicallyRealistic}. Delphi-2M and other large trajectory models frame this task as learning disease natural history, while digital-twin language emphasizes patient-level continuation \citep{Makarov2025Largelanguagemodels,Shmatko2025Delphi2M}. Such language should remain evidence-calibrated: generated continuity is not proof that a synthetic patient responds correctly to a changed intervention.

Two recent directions sharpen the boundary. One argues that a patient should be modeled as a changing latent system rather than as a moving document \citep{Adam2026Patientisnot}; another generates alternative patient timelines or ophthalmic trajectories \citep{Akagi2026GeneratingCounterfactualPatient,Belligoli2026OphthadtGenerativeDigital}. These are relevant to world modeling because they challenge static record prediction. Yet the word ``counterfactual'' in a title or interface does not establish Level~4: the alternative must be tied to a defined intervention and a defensible identification strategy. At Level~2, these models are best read as structured generators of possible futures.

Across this literature, outputs mix at least two dynamics: the patient's evolving condition and the healthcare system's response to it. A medication token may be generated because the model predicts clinician behavior rather than because the medication has been set as an intervention. Whether that distinction is resolved determines whether the system can be interpreted as a patient simulator.

\subsubsection{Physiological and Mechanistic Forecasting}

Physiological forecasting and digital-twin components also appear at Level~2. Graph-based patient-state forecasting, patient-specific hybrid simulation, aneurysm-flow surrogates, cardiac predictive representations, learned electrophysiology, and population virtual-patient databases each encode a temporal or mechanistic transition \citep{Barbiero2021GraphRepresentationForecasting,Sharma2025PatientSpecificDigital,Garnier2025GraphDeepLearning,Fernandes2026BeyondPatientInvariance,Zhou2026LearningCardiacElectrophysiology,Weber2026SyntheticPopulationBased}. A patient-specific autoregressive heart model similarly forecasts near-term wearable measurements \citep{Wickramasinghe2025HeartTwin}. The presence of equations, agents, or a digital-twin label does not determine the level. The relevant question is whether the audited experiment forecasts natural evolution, simulates a settable intervention, or uses the model to choose one.

These systems can support monitoring, early warning, cohort construction, and parameter estimation, but their temporal horizon and observation process must match the proposed use. Forecasting the next minute of a waveform, the next clinical event, and a five-year disease trajectory are not interchangeable forms of evidence. Mechanistic structure can improve extrapolation, yet its credibility depends on parameter identifiability and on validation against the physiological variables that the downstream decision actually uses.

\subsubsection{Surgical Video and Action-Labeled Boundary Cases}

Surgical video generation exposes this boundary directly. SurgSora, VISAGE, HieraSurg, and suturing world-model studies predict or generate procedural futures using object, hierarchy, or action annotations \citep{Chen2024SurgsoraObjectAware,Yeganeh2024VisageVideoSynthesis,Biagini2025HierasurgHierarchyAware,Turkcan2025SuturingWorldModels}. Expert assessment of zero-shot surgical video generation asks whether general video models capture clinically meaningful procedural dynamics \citep{Chen2025HowFarAre}. These studies are more informative than unconditional video synthesis, but an action label is not sufficient: the model must permit a feasible action to be changed for the same state and must validate the corresponding post-action transition before it enters Level~3.

Decision-learning comparisons also help define this boundary. DARIL, for example, contrasts imitation and reinforcement learning in surgical action planning \citep{Boels2025DarilWhenImitation}. Such work belongs in the world-model evidence map when it tests whether a predictive model contributes to action selection, but the policy result should not be attributed to a world model unless that transition model is explicit and ablated. We therefore retain these studies as boundary evidence rather than allowing planning terminology to determine their capability level.

Across applications, the common property is temporal direction without a validated, exposed decision variable. The models can support prognosis, missing-observation synthesis, trial enrichment, educational visualization, and anomaly detection. They cannot support treatment comparison merely because treatment appears somewhere in the history or because multiple futures can be sampled.

\subsection{Representative Model Designs}

Autoregressive transformers and language-model-like systems are natural for coded EHR events. They model a distribution over the next token or event sequence and can roll forward for long horizons. Their main representational challenge is to distinguish patient state from the process that records it. Recurrent networks and latent state-space models compress irregular histories and can expose an explicit latent transition. Neural ODEs and continuous-time latent models handle sparse or uneven intervals by defining evolution between observations.

Diffusion, flow, and adversarial models dominate future-image synthesis. They can preserve spatial detail and represent multimodal futures, which is valuable when disease trajectories are uncertain. Their output quality can nevertheless encourage an interpretation stronger than the training target. A visually convincing image may reflect learned population appearance rather than patient-specific biology. Mechanistic and hybrid designs provide a different bias: known equations constrain progression, while learned components estimate patient parameters or residual dynamics. These models may trade visual flexibility for interpretability and extrapolation structure.

Architecture should be selected according to the temporal question. Short-horizon observation forecasting may favor high-capacity sequence models; sparse longitudinal disease modeling may benefit from continuous-time or mechanistic constraints; image synthesis may require generative uncertainty; multimodal trajectories may require explicit alignment across measurement schedules. No family automatically resolves the action boundary. The same autoregressive or diffusion backbone can move to Level~3 only when a settable action changes the modeled future.

\subsection{Evidence Quality and Failure Modes}

Validation at Level~2 should be horizon-specific and should separate interpolation from extrapolation. One-step accuracy does not establish multi-step stability. Image studies should examine anatomy, pathology, identity preservation, and uncertainty in addition to perceptual quality. EHR studies should report patient-level temporal splits, event calibration, long-horizon degradation, missingness sensitivity, and performance under site or coding shift. Physiological models should evaluate whether errors remain acceptable at the timescale relevant to monitoring or intervention.

The most characteristic failure is to confuse prediction of the record with prediction of patient dynamics. Tests are ordered because a patient is deteriorating; medications are started in response to severity; discharge codes summarize a complex process. A model that reproduces those sequences may be useful for workflow forecasting while remaining unreliable for biological progression. A second failure is regression toward common trajectories, which can produce plausible average futures while missing rare deterioration. A third is horizon mismatch, in which a model validated at one step is discussed as a long-term simulator. A fourth is leakage through future-derived labels, preprocessing, or noncausal visit construction.

Temporal uncertainty deserves particular attention. The set of plausible futures should widen with horizon and distribution shift. Sampling diversity alone is not calibrated uncertainty, and a model can generate many images or event sequences that are all wrong in the same way. Evaluation should test coverage, calibration, and clinically important modes of failure. External cohorts with different observation and care processes are especially informative because they reveal whether the model learned patient evolution or local documentation practices.

\subsection{Level Takeaway and Next Threshold}

Level~2 is the point at which a model begins to represent dynamics rather than only state. It is already useful for prognosis and longitudinal reasoning. Its limitation is not simply the absence of a treatment input; it is that the predicted future is entangled with the policy and observation process that generated the data. Clear language about this distinction improves both the science and the clinical interpretation.

\levelsummary
{Future-image, EHR-trajectory, disease-progression, and physiological models demonstrate that clinically meaningful temporal prediction is feasible across several modalities and horizons.}
{Predicting the next clinical record does not necessarily mean predicting patient dynamics. Many apparently dynamic models learn a mixture of biology, observation, and care-process behavior.}
{To reach Level~3, a model must expose an executable action for the same pre-action state, use it in the transition, and validate a meaningful post-action future against time-only and no-action alternatives.}
